\section{总结：Selective computation 的两条路}

前面分别从 attention 状态选择和 expert 参数选择展开，现在把两条路线统一为 selective computation：不是所有 token 都需要看所有历史，也不是所有 token 都需要激活所有参数。真正的工程判断是节省的 dense 计算能否覆盖新增的路由、状态维护、负载不均与跨设备通信成本。


\begin{figure}[H]
\centering
\includegraphics[width=0.88\textwidth]{slide-059.jpg}
\caption{Slide 60: MoE summary}
\figsource{CS336 Spring 2026 Lecture 4 官方幻灯片。}
\end{figure}

\noindent\textbf{展开说明：}MoEs take advantage of sparsity: not all inputs need the full model.

\subsection{本章小结}
本节的共同线索是 selective computation：通过结构选择让 token 不必总是访问所有历史或所有参数。但选择性越强，越需要额外机制保证信息流、负载均衡和训练稳定。


